\documentclass[11pt,a4paper]{article}
\usepackage[T1]{fontenc}
\usepackage[utf8]{inputenc}
\usepackage{amsmath,amssymb,amsthm}
\usepackage[margin=1in]{geometry}
\usepackage[colorlinks=true,linkcolor=blue,urlcolor=blue]{hyperref}
\theoremstyle{plain}
\newtheorem{theorem}{Theorem}
\newtheorem{lemma}[theorem]{Lemma}
\newtheorem{proposition}[theorem]{Proposition}
\newtheorem{corollary}[theorem]{Corollary}
\theoremstyle{definition}
\newtheorem{remark}[theorem]{Remark}

\title{Recovering the unwritten recurrences of the table OEIS A237175}
\author{Adrian Perez Fontelles\\ \small Independent researcher}
\date{2 September 2026}
\begin{document}
\maketitle

\begin{abstract}
OEIS A237175 is a two-dimensional table $T(n,k)$ whose empirical recurrences are, for several of
its lines, given only as an ORDER --- ``k=2: [order 75]'' and the like --- with the
coefficients in a linked file rather than in the entry. Those conjectures can still be settled,
and the recurrences recovered. A line of the table is a fixed-width or fixed-height array
count, hence a walk count on finitely many vertices, hence satisfies some monic recurrence of
order at most that number; Berlekamp--Massey on twice as many exact terms returns the MINIMAL
one. Where its order equals the order the entry states --- and where the same holds after the
first few terms are dropped --- any recurrence of that order the line satisfies has a
characteristic polynomial that is a multiple of the minimal one and of equal degree, hence is
that polynomial. This note settles column 2.
\end{abstract}

\noindent\small 2020 Mathematics Subject Classification. 05A15, 11B37, 15A18.\normalsize

\section{The table and the unwritten conjectures}

OEIS A237175 is ``T(n,k)=Number of (n+1)X(k+1) 0..3 arrays with the maximum minus the lower median of every 2X2 subblock equal.''. It is read by antidiagonals and begins
\[
256, 1696, 1696, 11242, 25384, 11242, 79688, 370670,\ \dots
\]
The entry carries, among its empirical claims:
\begin{quote}\small
k=2: [order 75]
\end{quote}
As of the ``Last modified'' line on the live entry (Jul 23 2025 09:56:16, revision 6) these are still
recorded as empirical, and the recurrences themselves are not in the entry text.

\section{Why a line of the table is a walk count}

Fix the index that the line fixes. The entry defines $T(n,k)$ as a number of arrays of a shape
determined by $n$ and $k$, subject to a condition local to a bounded window of consecutive
lines. Substituting the fixed index into the entry's own wording turns the definition into an
ordinary one-parameter array count, and for such a count the admissible windows are the
vertices of a finite digraph and an array is a walk. So the line is a walk count on $S$
vertices, and by the Cayley--Hamilton theorem it satisfies a monic integer recurrence of order
at most $S$.

\section{Recovering the recurrences}

\begin{lemma}\label{lem:bm}
A sequence known to satisfy some linear recurrence of order at most $S$ is determined, with its
minimal recurrence, by its first $2S$ terms; Berlekamp--Massey applied to them returns that
minimal recurrence exactly.
\end{lemma}

\subsection*{Column $k=2$}
Substituting the fixed index into the entry's wording gives an ordinary one-parameter array
count, a walk on $S=202$ vertices. Berlekamp--Massey on $2S=404$ exact terms
returns a minimal recurrence of order $75$ --- the order the entry states --- with
integer coefficients
\[
(c_1,\dots,c_{75})=(42,\ -200,\ -13865,\ 178231,\ 1582899,\ -36847314,\ -30022430,\ 3930027822,\ -10796333296,\ -254160338775,\ 1351226532795,\ 10423480254120,\ -85854647426135,\ -256266292535981,\ 3562189924627891,\ 2181601254733383,\ -104444678293908555,\ 96643291574413178,\ 2236393294416712573,\ -4821647012167932615,\ -35230672810770863020,\ 119815925611797339580,\ 399302765019585986967,\ -2034763208847077077320,\ -2952981664132399365163,\ 25557714831009194351010,\ 7597196816489732891319,\ -244421424272298925917771,\ 133423378126754223121142,\ 1792797617824200067535664,\ -2285976172502440423434913,\ -9977303805235283369948651,\ 20469029382758672005780330,\ 40437055890040082135829545,\ -127693695090293973040308166,\ -104040805612161854388749529,\ 593392715973637567424597362,\ 49460050876741302982881651,\ -2091156696388820069165086299,\ 986150282278397970470985215,\ 5550957861367845992587119606,\ -5465127455770288417999745966,\ -10711476514584112231842228080,\ 17224727992861766578089894636,\ 13454189450330305517181041764,\ -37581411171840085624805802584,\ -5779357774683243988393019520,\ 59049871099596272036424832608,\ -16596664039327904452543332384,\ -66087286624193446947493541376,\ 44218230907086633175529522368,\ 49355354191342837407697083776,\ -57773689016050705899461247488,\ -19076059128455746505837351424,\ 47908113689601352879978061824,\ -3860759762679173282170552320,\ -25656404994628140323801157632,\ 9912358933854910371979370496,\ 8082984399737224234436460544,\ -6057437882095035177278210048,\ -911882703554153817390645248,\ 1905089366380530442491527168,\ -275682488441407650461122560,\ -297705684318409825173635072,\ 110544100184095301080973312,\ 13606945677734603130929152,\ -13208598734144929976549376,\ 1270874474326660638310400,\ 487357397477307958951936,\ -95251143522904361140224,\ -5164613546983336968192,\ 2024241162587908079616,\ -31052517093704466432,\ -13902362804150599680,\ 562082713624903680).
\]
so that $a(m)=\sum_{i=1}^{75}c_i\,a(m-i)$ for every $m$. Removing the first
$75+4$ terms and repeating gives order $75$ again, which is what the
identification below needs.

\begin{theorem}
For each line above, the displayed recurrence holds for every index, and it is the recurrence
the entry records.
\end{theorem}

\begin{proof}
It holds by Lemma~\ref{lem:bm}. For the identification, the entry asserts that the line
satisfies SOME recurrence of the stated order, possibly only from some index on. Let $q$ be its
characteristic polynomial and $q_{\min}$ the minimal polynomial of the tail. Then
$q_{\min}\mid q$, and the second Berlekamp--Massey run --- on the line with its first few
terms removed --- shows $\deg q_{\min}$ equals the stated order, which is $\deg q$. Both are
monic, so $q=q_{\min}$.
\end{proof}

\section{Verification}

Three checks.

First, each line's model was compared against the entry's own published values for that line,
taken from the antidiagonal DATA read in either orientation and from the ``Table starts''
block, and agrees at every available term. This is the only point at which reading the entry's
English enters, and a misreading gives different counts.

Second, each recovered recurrence was evaluated directly on the model's values, with no
Berlekamp--Massey involved, and holds at every index tested.

Third, each order was computed twice by independent routes --- modulo a $61$-bit prime, where
every intermediate is a machine word, and again in exact rational arithmetic --- and once more
on the line with its first few terms dropped, which is what the identification requires. All
agree.

\begin{thebibliography}{9}
\bibitem{oeis} The OEIS Foundation, \emph{The On-Line Encyclopedia of Integer
Sequences}, \url{https://oeis.org}, 2026. Sequence A237175.
\bibitem{massey} J.~L.~Massey, Shift-register synthesis and BCH decoding, \emph{IEEE
Trans. Inform. Theory} \textbf{15} (1969), 122--127.
\bibitem{stanley} R.~P.~Stanley, \emph{Enumerative Combinatorics}, Volume 1, 2nd ed.,
Cambridge University Press, 2011. (Transfer-matrix method, Section 4.7.)
\bibitem{horn} R.~A.~Horn and C.~R.~Johnson, \emph{Matrix Analysis}, 2nd ed., Cambridge
University Press, 2013.
\end{thebibliography}

\end{document}
